\documentclass[11pt]{article}

\usepackage[margin=1in]{geometry}
\usepackage{amsmath}
\usepackage{amssymb}
\usepackage{xcolor}
\usepackage[most]{tcolorbox}
\usepackage{enumitem}
\usepackage{titlesec}
\usepackage{hyperref}

% Links map to the structural header blue.
% bookmarks=false avoids the spaced-jobname .out failure on Windows/tectonic.
\hypersetup{colorlinks=true, linkcolor=headerblue, urlcolor=headerblue, bookmarks=false}

\tcbuselibrary{skins}

% 1.1 Palette
\definecolor{headerblue}{HTML}{1C569E}
\definecolor{accentgreen}{HTML}{2E7D32}
\definecolor{accentorange}{HTML}{E27A1E}
\definecolor{workpurple}{HTML}{A020F0}
\definecolor{warnred}{HTML}{C0392B}

% 1.2 Subsection Typography
\titleformat{\subsection}
  {\normalfont\large\bfseries\color{headerblue}}{\thesubsection}{1em}{}

% 2. Table of Contents Depth
\setcounter{secnumdepth}{2}
\setcounter{tocdepth}{2}

% 3. Box System (Unbreakable)
\newtcolorbox{topicsbox}{
    enhanced,
    colback=headerblue!6!white, colframe=headerblue,
    coltitle=white, fonttitle=\bfseries,
    title={Convolution at a Glance: Core Sub-Topics},
    boxrule=0.6pt, arc=2mm,
    left=3mm, right=3mm, top=2mm, bottom=2mm
}

\newtcolorbox{formulabox}[1][Master Formula]{
    enhanced,
    colback=accentorange!8!white, colframe=accentorange,
    coltitle=white, fonttitle=\bfseries,
    title={#1},
    boxrule=0.6pt, arc=1mm,
    left=3mm, right=3mm, top=2mm, bottom=2mm
}

\newtcolorbox{methodbox}[1][Method]{
    enhanced,
    colback=accentgreen!6!white, colframe=accentgreen,
    coltitle=white, fonttitle=\bfseries,
    title={#1},
    boxrule=0.6pt, arc=1mm,
    left=3mm, right=3mm, top=2mm, bottom=2mm
}

\newtcolorbox{examplebox}[1][Worked Example]{
    enhanced,
    colback=workpurple!4!white, colframe=workpurple,
    coltitle=white, fonttitle=\bfseries,
    title={#1},
    boxrule=0.4pt, sharp corners,
    borderline west={3pt}{0pt}{workpurple},
    left=5mm, right=3mm, top=2mm, bottom=2mm
}

\newtcolorbox{conceptbox}[1][Key Takeaway]{
    enhanced,
    colback=accentorange!5!white, colframe=accentorange,
    coltitle=white, fonttitle=\bfseries,
    title={#1},
    boxrule=0.3pt, arc=1mm,
    left=3mm, right=3mm, top=2mm, bottom=2mm
}

\newtcolorbox{warningbox}[1][Common Pitfall]{
    enhanced,
    colback=red!3!white, colframe=warnred,
    coltitle=white, fonttitle=\bfseries,
    title={#1},
    boxrule=0.4pt, sharp corners,
    borderline west={3pt}{0pt}{warnred},
    left=5mm, right=3mm, top=2mm, bottom=2mm
}

\begin{document}

\begin{center}
    {\Large\bfseries\color{headerblue} Convolution, Duhamel's Principle,\\[2pt] and Volterra Equations}\\[6pt]
    \large A Self-Study Reference \& Master Guide
\end{center}
\vspace{0.4cm}

\tableofcontents
\vspace{0.8cm}

%======================================================================
\phantomsection
\addcontentsline{toc}{section}{Part I --- Theory \& Worked Examples}
\section*{Part I --- Theory \& Worked Examples}
\setcounter{section}{4}
\setcounter{subsection}{1}

\begin{topicsbox}
\begin{itemize}[leftmargin=*, itemsep=2pt]
  \item \textbf{Algebraic Properties (4.2):} convolution as a true ``product'' in the time domain --- commutative, associative, distributive, scalar-linear, and annihilated by zero --- plus the pitfalls that separate it from ordinary multiplication.
  \item \textbf{Duhamel's Principle (4.3):} the impulse-response viewpoint. Once the response $h(t)$ to a unit impulse is known, the system's response to \emph{any} input is the convolution $h * f$.
  \item \textbf{Volterra Integral Equations (4.4):} equations with \emph{memory}, where the unknown $x(t)$ appears inside a convolution integral. The transform turns them into one-line algebra in $s$.
\end{itemize}
\end{topicsbox}

\subsection{Convolution --- Properties and Pitfalls}

The convolution theorem tells us that a product of transforms $F(s)G(s)$ inverts to the convolution $(f * g)(t) = \int_0^t f(\tau)\,g(t-\tau)\,d\tau$. Before we put this to work, it pays to treat convolution as an algebraic operation in its own right --- one that behaves like multiplication in some ways, and dangerously \emph{unlike} it in others.

\begin{conceptbox}[Algebraic Properties]
For piecewise-continuous, exponentially-bounded functions $f, g, h$ and any constant $c$, convolution obeys the familiar algebra of a product:
\begin{itemize}[leftmargin=*, itemsep=1pt]
  \item \textcolor{accentorange}{Commutative:} $f * g = g * f$.
  \item \textcolor{accentorange}{Associative:} $f * (g * h) = (f * g) * h$.
  \item \textcolor{accentorange}{Distributive:} $f * (g + h) = f * g + f * h$.
  \item \textcolor{accentorange}{Scalar-linear:} $(c f) * g = c\,(f * g)$.
  \item \textcolor{accentorange}{Annihilated by zero:} $0 * f = 0$.
\end{itemize}
Each one mirrors a corresponding fact about the product $F(s)\,G(s)$ in the $s$-domain, which is why they hold.
\end{conceptbox}

\begin{warningbox}[No Multiplicative Identity; $f*f$ Can Be Negative]
Convolution looks like multiplication, but two instincts carried over from ordinary products will burn you.

\smallskip
\textbf{1. The constant $1$ is \emph{not} the identity: $f * 1 \ne f$.} Convolving with $1$ integrates:
\[ (f * 1)(t) = \int_0^t f(\tau)\cdot 1\, d\tau = \int_0^t f(\tau)\, d\tau. \]
For example $(\cos t) * 1 = \int_0^t \cos\tau\,d\tau = \sin t \ne \cos t$. The true convolution identity is the \textbf{Dirac delta}: by the sifting property,
\[ (f * \delta)(t) = \int_0^t f(\tau)\,\delta(t-\tau)\,d\tau = f(t). \]

\smallskip
\textbf{2. $f * f$ is not a square.} Unlike $f(t)^2$, which is never negative, the convolution $f * f$ can dip below zero. For instance $(\sin t)*(\sin t) = \tfrac{1}{2}(\sin t - t\cos t)$, which is negative on a range of $t$. Never assume $f * f \ge 0$.
\end{warningbox}

\begin{examplebox}[Worked Example 1: $(\cos t) * (\cos t)$]
We evaluate the self-convolution of cosine directly from the definition.

\smallskip
\textbf{Step 1 --- Write the convolution integral.}
\[ \color{workpurple} (\cos * \cos)(t) = \int_0^t \cos(t-\tau)\,\cos\tau\,d\tau. \]

\textbf{Step 2 --- Apply the product-to-sum identity.} Using $\cos A\cos B = \tfrac{1}{2}\big[\cos(A-B) + \cos(A+B)\big]$ with $A = t-\tau$ and $B = \tau$, the difference is $A-B = t - 2\tau$ and the sum is $A+B = t$:
\[ \color{workpurple} \cos(t-\tau)\cos\tau = \tfrac{1}{2}\big[\cos(t-2\tau) + \cos t\big]. \]
So the integral becomes
\[ \color{workpurple} (\cos * \cos)(t) = \frac{1}{2}\int_0^t \big[\cos(t-2\tau) + \cos t\big]\,d\tau. \]

\textbf{Step 3 --- Integrate term by term in $\tau$.} The variable of integration is $\tau$, so the standalone $\cos t$ is a \emph{constant} here --- it pulls straight through the integral and contributes a factor of $t$:
\[ \color{workpurple} \int_0^t \cos t\,d\tau = t\cos t. \]
For the other term, substitute $u = t - 2\tau$, $du = -2\,d\tau$, so $\int \cos(t-2\tau)\,d\tau = -\tfrac{1}{2}\sin(t-2\tau)$. Evaluating from $\tau = 0$ to $\tau = t$:
\[ \color{workpurple} \int_0^t \cos(t-2\tau)\,d\tau = \Big[-\tfrac{1}{2}\sin(t-2\tau)\Big]_0^t = -\tfrac{1}{2}\sin(-t) + \tfrac{1}{2}\sin(t) = \sin t. \]

\textbf{Step 4 --- Combine.}
\[ \color{workpurple} (\cos * \cos)(t) = \tfrac{1}{2}\big[\sin t + t\cos t\big]. \]
The linearly growing factor $t\cos t$ is the fingerprint of resonance --- a preview of Section 4.3.
\end{examplebox}

\subsection{Duhamel's Principle --- Response to Arbitrary Input}

Classical methods solve $x'' + \omega_0^2 x = f(t)$ one forcing function at a time. Duhamel's principle does it \emph{once and for all}: it produces a single integral formula valid for \emph{every} input $f$, built entirely from the system's response to a unit impulse.

\begin{methodbox}[The Killer Application --- Four Steps]
Solve $x'' + \omega_0^2\, x = f(t)$ with rest initial conditions $x(0) = x'(0) = 0$, for an \emph{arbitrary} forcing $f(t)$.

\smallskip
\textbf{Step 1 --- Transform.} With zero initial conditions, $\mathcal{L}\{x''\} = s^2 X(s)$, so
\[ s^2 X + \omega_0^2 X = F(s) \quad\Longrightarrow\quad X(s) = \frac{F(s)}{s^2 + \omega_0^2}. \]

\textbf{Step 2 --- Identify the transfer function and impulse response.} Factor off the input:
\[ X(s) = H(s)\,F(s), \qquad H(s) = \frac{1}{s^2 + \omega_0^2}, \qquad h(t) = \mathcal{L}^{-1}\{H(s)\} = \frac{\sin(\omega_0 t)}{\omega_0}. \]
Here $h(t)$ is the \emph{impulse response} --- the output when $f = \delta(t)$, since $\mathcal{L}\{\delta\} = 1$.

\textbf{Step 3 --- Convolve.} A product in $s$ is a convolution in $t$, so
\[ x(t) = (h * f)(t) = \int_0^t \frac{\sin\big(\omega_0(t-\tau)\big)}{\omega_0}\, f(\tau)\,d\tau. \]
This single integral solves the ODE for \emph{any} admissible $f$.

\textbf{Step 4 --- Resonance check.} Feed the natural frequency back in: $f(t) = \cos(\omega_0 t)$. The convolution then produces a linearly growing amplitude,
\[ x(t) = \frac{t\,\sin(\omega_0 t)}{2\omega_0}, \]
the unbounded growth that signals resonance. The full computation is expanded below.
\end{methodbox}

\begin{examplebox}[Expanding Step 4: Driving at the Natural Frequency]
We carry out the convolution of Step 3 with $f(\tau) = \cos(\omega_0\tau)$ in full.

\smallskip
\textbf{Step 1 --- Set up the integral.}
\[ \color{workpurple} x(t) = \int_0^t \frac{\sin\big(\omega_0(t-\tau)\big)}{\omega_0}\,\cos(\omega_0\tau)\,d\tau = \frac{1}{\omega_0}\int_0^t \sin\big(\omega_0(t-\tau)\big)\cos(\omega_0\tau)\,d\tau. \]

\textbf{Step 2 --- Product-to-sum.} Use $\sin A\cos B = \tfrac{1}{2}\big[\sin(A+B) + \sin(A-B)\big]$ with $A = \omega_0(t-\tau)$ and $B = \omega_0\tau$. Then $A+B = \omega_0 t$ and $A-B = \omega_0(t - 2\tau)$:
\[ \color{workpurple} x(t) = \frac{1}{2\omega_0}\int_0^t \Big[\sin(\omega_0 t) + \sin\big(\omega_0(t-2\tau)\big)\Big]\,d\tau. \]

\textbf{Step 3 --- Integrate in $\tau$.} The term $\sin(\omega_0 t)$ is constant in $\tau$, so it integrates to $t\sin(\omega_0 t)$. The second term integrates to zero over $[0,t]$:
\[ \color{workpurple} \int_0^t \sin\big(\omega_0(t-2\tau)\big)\,d\tau = \Big[\tfrac{1}{2\omega_0}\cos\big(\omega_0(t-2\tau)\big)\Big]_0^t = \tfrac{1}{2\omega_0}\big[\cos(-\omega_0 t) - \cos(\omega_0 t)\big] = 0, \]
because $\cos(-\theta) = \cos\theta$ makes the two boundary terms cancel.

\textbf{Step 4 --- Combine.}
\[ \color{workpurple} x(t) = \frac{1}{2\omega_0}\,t\,\sin(\omega_0 t) = \frac{t\,\sin(\omega_0 t)}{2\omega_0}. \]
The amplitude grows without bound as $t \to \infty$: this is resonance, falling out of the convolution automatically when the forcing frequency matches the natural frequency $\omega_0$.
\end{examplebox}

\begin{conceptbox}[Key Takeaway]
Once the \textcolor{accentorange}{impulse response $h(t)$} is known, the system's response to \emph{any} input is the convolution $x = h * f$. \textcolor{accentorange}{The system is fully characterized by $h$} --- no new solve is needed when the forcing changes.
\end{conceptbox}

\subsection{Volterra Integral Equations}

Differential equations describe systems whose future depends only on their \emph{present} state. Many physical systems instead carry \emph{memory}: the present response depends on the entire history of the input. Viscoelastic materials ``remember'' past deformation; age-structured populations depend on their full birth history; radiative transfer accumulates contributions along a path. These are modeled by \textbf{integral equations}, where the unknown sits inside a convolution. The Laplace transform dissolves the memory integral into ordinary algebra.

\begin{formulabox}[Solving by Algebra in $s$]
A Volterra integral equation of convolution type has the form
\[ x(t) = f(t) + \int_0^t g(t-\tau)\,x(\tau)\,d\tau. \]
The integral is exactly the convolution $(g * x)(t)$, so applying $\mathcal{L}$ and using $\mathcal{L}\{g * x\} = G(s)X(s)$:
\[ X(s) = F(s) + G(s)\,X(s) \quad\Longrightarrow\quad X(s)\big[1 - G(s)\big] = F(s) \quad\Longrightarrow\quad \boxed{\,X(s) = \frac{F(s)}{1 - G(s)}\,}. \]
The unknown is isolated by simple division --- the integral equation has become an algebraic one.
\end{formulabox}

\begin{examplebox}[Worked Example 2: A Volterra Equation with a $\sinh$ Kernel]
Solve $\displaystyle x(t) = e^{-t} + \int_0^t \sinh(t-\tau)\,x(\tau)\,d\tau$.

\smallskip
\textbf{Step 1 --- Transform both sides.} The integral is the convolution $(\sinh * x)(t)$. With $\mathcal{L}\{e^{-t}\} = \tfrac{1}{s+1}$ and $\mathcal{L}\{\sinh t\} = \tfrac{1}{s^2 - 1}$:
\[ \color{workpurple} X(s) = \frac{1}{s+1} + \frac{1}{s^2 - 1}\,X(s). \]

\textbf{Step 2 --- Collect the $X$ terms.} Move the $X$ term to the left and factor it out:
\[ \color{workpurple} X(s) - \frac{1}{s^2-1}\,X(s) = \frac{1}{s+1} \quad\Longrightarrow\quad X(s)\left(1 - \frac{1}{s^2-1}\right) = \frac{1}{s+1}. \]

\textbf{Step 3 --- Common denominator inside the bracket.}
\[ \color{workpurple} 1 - \frac{1}{s^2-1} = \frac{(s^2-1) - 1}{s^2-1} = \frac{s^2 - 2}{s^2 - 1}, \qquad\text{so}\qquad X(s)\cdot\frac{s^2-2}{s^2-1} = \frac{1}{s+1}. \]

\textbf{Step 4 --- Multiply by the reciprocal and cancel.} Solving for $X(s)$ and using $s^2 - 1 = (s-1)(s+1)$:
\[ \color{workpurple} X(s) = \frac{1}{s+1}\cdot\frac{s^2-1}{s^2-2} = \frac{(s-1)(s+1)}{(s+1)(s^2-2)} = \frac{s-1}{s^2-2}. \]

\textbf{Step 5 --- Split into table-ready pieces.} Decompose by separating the numerator:
\[ \color{workpurple} X(s) = \frac{s-1}{s^2-2} = \frac{s}{s^2-2} - \frac{1}{s^2-2} = \frac{s}{s^2-2} - \frac{1}{\sqrt{2}}\cdot\frac{\sqrt{2}}{s^2-2}. \]

\textbf{Step 6 --- Invert.} Matching $\tfrac{s}{s^2-a^2} \to \cosh(at)$ and $\tfrac{a}{s^2-a^2} \to \sinh(at)$ with $a = \sqrt{2}$:
\[ \color{workpurple} x(t) = \cosh\!\big(\sqrt{2}\,t\big) - \frac{1}{\sqrt{2}}\,\sinh\!\big(\sqrt{2}\,t\big). \]
\end{examplebox}

\begin{conceptbox}[Key Takeaway]
A convolution-type integral equation transforms to $X(s) = \dfrac{F(s)}{1 - G(s)}$ --- the \textcolor{accentorange}{memory integral becomes a single division}. Integral equations of physics-with-memory (viscoelasticity, age-structured populations, radiative transfer) yield to Laplace just as cleanly as differential equations.
\end{conceptbox}

%======================================================================
\clearpage
\phantomsection
\addcontentsline{toc}{section}{Part II --- Practice Set}
\section*{Part II --- Practice Set}

Work each problem before turning to the complete solutions in Part III. Both reward setting up the integral or the $s$-domain algebra carefully before computing.

\bigskip
\noindent\textbf{Problem 1: A Convolution Evaluation.}\\[2pt]
Evaluate the convolution $t * \sin t$ directly from the definition,
\[ (t * \sin t)(t) = \int_0^t (t-\tau)\,\sin\tau\,d\tau, \]
and confirm your answer using the convolution theorem $\mathcal{L}\{t * \sin t\} = \mathcal{L}\{t\}\cdot\mathcal{L}\{\sin t\}$.

\bigskip
\noindent\textbf{Problem 2: A Volterra Integral Equation.}\\[2pt]
Solve for $x(t)$:
\[ x(t) = t + \int_0^t \sin(t-\tau)\,x(\tau)\,d\tau. \]
\textit{Hint: the integral is a convolution $(\sin * x)(t)$; transform, isolate $X(s)$, and invert.}

%======================================================================
\clearpage
\phantomsection
\addcontentsline{toc}{section}{Part III --- Complete Solutions}
\section*{Part III --- Complete Solutions}

\noindent\textbf{Solution 1: $t * \sin t$.}\\[2pt]
We evaluate the integral directly, then verify in the $s$-domain.

\smallskip
\textbf{Step 1 --- Split the integral.} Distribute $(t - \tau)$, remembering that $t$ is constant with respect to the integration variable $\tau$:
\[ \color{workpurple} \int_0^t (t-\tau)\sin\tau\,d\tau = t\int_0^t \sin\tau\,d\tau - \int_0^t \tau\sin\tau\,d\tau. \]

\noindent\textbf{Step 2 --- First piece.}
\[ \color{workpurple} t\int_0^t \sin\tau\,d\tau = t\big[-\cos\tau\big]_0^t = t\,(1 - \cos t) = t - t\cos t. \]

\noindent\textbf{Step 3 --- Second piece by parts.} Take $u = \tau$, $dv = \sin\tau\,d\tau$, so $du = d\tau$, $v = -\cos\tau$:
\[ \color{workpurple} \int_0^t \tau\sin\tau\,d\tau = \big[-\tau\cos\tau\big]_0^t + \int_0^t \cos\tau\,d\tau = -t\cos t + \sin t. \]

\noindent\textbf{Step 4 --- Combine.}
\[ \color{workpurple} (t * \sin t)(t) = (t - t\cos t) - (-t\cos t + \sin t) = t - \sin t. \]

\noindent\textbf{Step 5 --- Verify via the convolution theorem.}
\[ \color{workpurple} \mathcal{L}\{t * \sin t\} = \frac{1}{s^2}\cdot\frac{1}{s^2+1} = \frac{1}{s^2(s^2+1)} = \frac{1}{s^2} - \frac{1}{s^2+1}, \]
which inverts term by term to $\textcolor{workpurple}{t - \sin t}$ --- in agreement with the direct integration.

\bigskip
\noindent\textbf{Solution 2: The Volterra Equation $x(t) = t + \int_0^t \sin(t-\tau)x(\tau)\,d\tau$.}\\[2pt]
The integral is the convolution $(\sin * x)(t)$, so we transform and solve for $X(s)$.

\smallskip
\textbf{Step 1 --- Transform.} With $\mathcal{L}\{t\} = \tfrac{1}{s^2}$ and $\mathcal{L}\{\sin t\} = \tfrac{1}{s^2+1}$:
\[ \color{workpurple} X(s) = \frac{1}{s^2} + \frac{1}{s^2+1}\,X(s). \]

\textbf{Step 2 --- Isolate $X(s)$.} Collect terms and factor:
\[ \color{workpurple} X(s)\left(1 - \frac{1}{s^2+1}\right) = \frac{1}{s^2}, \qquad 1 - \frac{1}{s^2+1} = \frac{(s^2+1) - 1}{s^2+1} = \frac{s^2}{s^2+1}. \]

\textbf{Step 3 --- Multiply by the reciprocal.}
\[ \color{workpurple} X(s)\cdot\frac{s^2}{s^2+1} = \frac{1}{s^2} \quad\Longrightarrow\quad X(s) = \frac{1}{s^2}\cdot\frac{s^2+1}{s^2} = \frac{s^2+1}{s^4}. \]

\textbf{Step 4 --- Split into table-ready pieces.}
\[ \color{workpurple} X(s) = \frac{s^2+1}{s^4} = \frac{1}{s^2} + \frac{1}{s^4}. \]

\textbf{Step 5 --- Invert.} Using $\mathcal{L}^{-1}\{1/s^2\} = t$ and $\mathcal{L}^{-1}\{1/s^4\} = \tfrac{t^3}{3!} = \tfrac{t^3}{6}$:
\[ \color{workpurple} x(t) = t + \frac{t^3}{6}. \]

\begin{conceptbox}[Closing Takeaway]
Both problems collapse to the same two moves: recognize the \textcolor{accentorange}{convolution structure}, then trade it for a \textcolor{accentorange}{product (or a division) in $s$}. Direct integration and the $s$-domain shortcut must always agree --- a built-in check on every convolution computation.
\end{conceptbox}

\end{document}
